% 심사용(저자 정보 제외)은 [preprint], 출판용은 옵션 없이 컴파일한다(KCC 템플릿 규칙).
\documentclass{kcc}
% ReHwp-Bench-Paper-copy/kcc.tex의 공통 프리앰블에서 계승.
\usepackage{graphicx}
\usepackage{amsmath}
\usepackage{amssymb}
\usepackage{array}
\usepackage{booktabs}
\usepackage{tabularx}
\usepackage{float}
\usepackage{url}
\usepackage{caption}
\usepackage{tikz}
\usetikzlibrary{arrows.meta,positioning,patterns,decorations.pathreplacing}
\usepackage{pgfplots}
\pgfplotsset{compat=1.18}
\usepackage{balance}
% \raggedbottom에서도 각주를 단 하단에 배치한다.
\usepackage[bottom]{footmisc}

\captionsetup{font=small,skip=3pt,labelsep=colon}
\setlength{\abovecaptionskip}{3pt}
\setlength{\belowcaptionskip}{3pt}
\clubpenalty=10000
\widowpenalty=10000
% 짧은 2단 원고의 빈 공간을 절·문단 사이로 늘리지 않는다.
\raggedbottom
\renewcommand{\tablename}{표}
\renewcommand{\figurename}{그림}
\renewcommand{\refname}{참고문헌}
\newcommand{\ph}[1]{\ensuremath{\langle #1\rangle}}
\newcommand{\draftnote}[1]{\par\noindent\textbf{[작성 메모]} #1\par}
\newcolumntype{L}{>{\raggedright\arraybackslash}X}
\newcommand{\papertablestyle}{\small\renewcommand{\arraystretch}{1.15}}

\tikzset{
  paperbox/.style={draw,rounded corners=1.2pt,align=center,font=\footnotesize,
    inner xsep=1.6pt,inner ysep=1.5pt,minimum height=5mm,minimum width=26mm},
  paperresult/.style={paperbox,fill=black!7},
  paperarrow/.style={-{Latex[length=1.4mm,width=1mm]},semithick},
  paperlabel/.style={font=\scriptsize,fill=white,inner sep=0.7pt}
}


% 참고문헌·시각자료를 포함해 4페이지 이내(실험 current, 현재 3페이지 목표).
% 용어: 편집 용이성(평가 관점), 양식(기입용 문서), 출력본(복원 입력·외관 기준),
% 원본(출력본을 생성한 워드프로세서 문서·편집 시간 기준), 복원본, 완성본.
% 원본은 수정 전 워드프로세서 문서이며, 출력본은 이를 출력한 PDF다.
% 글꼴·간격·정렬은 글자·문단 모양, 셀·병합은 문서 구조, 위치는 배치로 구분한다.
\title{PDF에서 복원한 워드프로세서 문서의 편집 시간 추정 기반 편집 용이성 평가}
\author{
정치우$^{\circ}$, 김민제, 하승민, 송하주\\
국립부경대학교 컴퓨터·인공지능공학부\\
cldnwjd1024@gmail.com, kmz34180@naver.com, ccy0111@naver.com, hajoosong@pknu.ac.kr
}
\engtitle{Evaluating the Ease of Editing Word Processor Documents Reconstructed from PDFs Using Estimated Editing Time}
\engauthor{
Chi-Woo Jung$^{\circ}$, Min-Je Kim, Seung-Min Ha, Ha-Joo Song\\
School of Computer Engineering \& Artificial Intelligence, Pukyong National University\\
cldnwjd1024@gmail.com, kmz34180@naver.com, ccy0111@naver.com, hajoosong@pknu.ac.kr
}
\abstract{
PDF로 배포된 양식을 워드프로세서에서 수정하려면 외관과 함께 문단과 표 등의 서식 요소를 복원할 필요가 있다.
그러나 복원본이 출력본과 비슷하게 보이더라도 내용을 입력하거나 표를 수정하는 데 많은 작업이 필요할 수 있다.
본 논문은 이러한 편집 용이성을 평가하기 위해 편집 시간을 추정하는 지표를 제안하고, 시각 유사도와 함께 측정한다.
시각 유사도는 복원본과 출력본의 외관을 비교하며, 편집 시간은 복원본을 완성본으로 만드는 데 필요한 시간을 추정하여 나타낸다.
원본 321개와 각각에 대응하는 완성본으로 네 가지 복원 방식을 평가한 결과, 편집 시간은 원본의 2.34--2.77배였다.
또한 누락되거나 다시 만들어야 하는 서식 요소를 추가하는 작업이 전체 추정 편집 시간의 54--64\%를 차지했다.
이 결과는 PDF 양식 복원 평가에서 시각 유사도와 함께 편집 용이성을 고려할 필요가 있음을 보여준다.

\vspace{2mm}
\noindent{\bfseries 키워드:} 양식 복원, 편집 용이성, 시각 유사도, 편집 시간, 평가 지표
}

\begin{document}
\maketitle
% 사사 각주는 출판본에만 표시한다.
\makeatletter
\if@preprint\else
\begingroup
\renewcommand{\thefootnote}{}
% 번호 없는 사사 각주는 내어쓰기 없이 정렬한다.
\long\def\@makefntext#1{\noindent#1}
\footnotetext{\fontsize{9}{11}\selectfont
본 연구는 2026년도 과학기술정보통신부 및 정보통신기획평가원의 \mbox{SW중심대학지원사업}의 연구결과로 수행되었음.}
\endgroup
\fi
\makeatother

\section{서 론}

PDF로 배포된 신청서를 워드프로세서에서 수정하려면, 이름이나 날짜를 입력하고 표의 행을 추가하는 등의 작업을 할 수 있어야 한다.
본 논문에서는 PDF 가져오기 기능이나 변환 도구로 PDF를 워드프로세서용 형식으로 바꾸는 과정을 \textbf{복원}이라 한다.
복원에는 외관의 재현뿐 아니라 이후 작업의 \textbf{편집 용이성}, 즉 내용을 기입하거나 구조를 고치는 데 얼마나 적은 작업이 필요한지도 중요하다.

이하 수정 전의 워드프로세서 문서를 \textbf{원본}, 이를 PDF로 출력한 결과를 \textbf{출력본}이라 한다.
출력본의 변환 결과를 \textbf{복원본}, 기입을 마친 목표 양식을 \textbf{완성본}이라 한다.

같은 값을 기입하더라도 복원본의 구조에 따라 필요한 작업이 달라질 수 있다.
그림~\ref{fig:cost-example}은 이러한 차이를 성명, 생년월일, 주소를 기입하는 가상 양식으로 설명한다.
원본에서는 셀에 값을 입력하는 데 약 24초, 행 하나가 누락된 복원본 A에서는 행 추가까지 약 31초로 추정된다.
표를 이미지로 삽입한 복원본 B는 출력본과 외관이 같더라도 표를 새로 만들어야 하므로 약 52초로 추정된다.
이처럼 외관이 유사해도 수정에 필요한 작업은 달라진다.
자세한 시간 추정 방법은 2.2절에서 설명한다.

{\setlength{\intextsep}{6pt plus 2pt minus 2pt}% 가상 신청서에 configs/costs_klm.yaml을 적용한 예. 수집 자료의 평가값이 아니다.
% src/formbench/editing/cost.py align(): FORMBENCH_COST_UNITS=klm, lang=ko, 허용오차 12pt.
% 목표: 성명=홍길동, 생년월일=1990. 1. 1., 주소=서울시 중구 (본문에서 설명).
% 원본 23.96초: 값 세 칸 입력. 행 누락 31.35초: 행 추가 + 주소 항목 + 값 세 칸 입력.
% 이미지 52.09초: 이미지 삭제 1.58 + 표 추가 50.51(표 삽입 8.65 + 여섯 칸 입력).
% 서론의 직관 예시(교수 피드백 2026-10-19): 외관의 유사성과 편집 시간이 어긋남을 보인다. 막대 높이는 반올림 전 편집 시간에 비례, 표시는 초.
% 서론 둘째 문단 바로 뒤에 오도록 [H]로 두고(사용자 요청 2026-10-20), 1페이지 왼쪽 단의 남은 높이에 맞게 낮게 그린다:
% 막대는 값 칸 아래에 두고, "편집 시간"은 표의 항목 칸 아래에 행 머리로 둔다.
\begin{figure}[H]
\centering
\begin{tikzpicture}[
  lc/.style={draw,minimum width=12mm,minimum height=3.8mm,inner sep=0pt,font=\scriptsize},
  vc/.style={draw,minimum width=15mm,minimum height=3.8mm,inner sep=0pt,font=\scriptsize},
  head/.style={font=\footnotesize,anchor=south,align=center,inner sep=1pt},
  value/.style={font=\footnotesize\bfseries,anchor=south,inner sep=1.5pt}]
\newcommand{\blankform}[1]{%
  \node[lc] at (#1+0.60,-0.19) {성명}; \node[vc] at (#1+1.95,-0.19) {};
  \node[lc] at (#1+0.60,-0.57) {생년월일}; \node[vc] at (#1+1.95,-0.57) {};}
\node[head] at (1.35,0.04) {원본\\{\scriptsize 기입만 필요}};
\node[head] at (4.55,0.04) {복원본 A\\{\scriptsize 행 하나 누락}};
\node[head] at (7.75,0.04) {복원본 B\\{\scriptsize 표를 이미지로 삽입}};
\blankform{0}
\node[lc] at (0.60,-0.95) {주소}; \node[vc] at (1.95,-0.95) {};
\blankform{3.2}
\draw[dashed,black!55] (3.2,-0.76) rectangle (5.9,-1.14) (4.4,-0.76) -- (4.4,-1.14);
\begin{scope}[every node/.append style={fill=black!15,draw=black!55,text=black!55}]
  \blankform{6.4}
  \node[lc] at (7.00,-0.95) {주소}; \node[vc] at (8.35,-0.95) {};
\end{scope}
% 공통의 0 기준선 위에서 길이만으로 시간 차이를 읽는다. 막대는 각 표의 값 칸 아래.
\draw[black!50] (1.30,-2.28) -- (9.10,-2.28);
\node[font=\footnotesize,anchor=south,inner sep=1pt] at (0.60,-2.28) {편집 시간};
\foreach \x/\ratio/\display in {1.95/1/24초,5.15/1.30843/31초,8.35/2.17404/52초} {
  \fill[black!60] (\x-0.40,-2.28) rectangle (\x+0.40,-2.28+0.32*\ratio);
  \node[value] at (\x,-2.28+0.32*\ratio) {\display};
}
\end{tikzpicture}
\caption{외관의 유사성과 편집 시간이 어긋나는 예}
\label{fig:cost-example}
\end{figure}
}

OmniDocBench는 PDF에서 추출한 텍스트, 표, 수식, 읽기 순서의 정확성을 평가하는 벤치마크를 제시했다~\cite{ouyang2025omnidocbench}.
PDF나 그림을 다시 활용할 수 있는 형식으로 복원하는 평가도 제안되었다.
TexOCR-Bench는 학술 PDF를 LaTeX로 복원한 결과의 내용, 구조, 컴파일 가능성을 평가하며~\cite{wang2026texocr}, ReFigBench는 학술 그림을 PowerPoint로 복원한 결과의 외관과 개체 구조를 평가한다~\cite{fan2026refigbench}.

이처럼 선행 연구는 정보 추출에서 복원 결과의 활용 가능성으로 평가 범위를 넓혀 왔다.
그러나 워드프로세서에서 양식에 기입하고 구조를 수정하는 데 필요한 작업은 이들 벤치마크의 평가 범위에 포함되지 않는다.

이에 본 논문은 복원본과 원본에서 같은 완성본을 만드는 \textbf{편집 시간}을 비교하는 평가 지표를 제안하고, 네 가지 복원 방식에 적용해 시각 유사도와 함께 분석한다.

\section{평가 지표}

그림~\ref{fig:pipeline}은 평가 과정을 나타낸다.
출력본을 \textbf{복원 방식}에 입력해 복원본을 생성하고, 시각 유사도와 편집 시간을 각각 계산한다.
시각 유사도의 비교 대상은 출력본이며, 편집 시간의 비교 대상은 원본에서 같은 완성본을 만드는 시간이다.
복원 방식에는 완성본과 원본을 제공하지 않고, 출력본만을 제공한다.

복원본의 형식은 국제 표준(ISO/IEC 26300)인 개방형 문서 형식 ODF의 텍스트 문서 형식 ODT로 통일한다~\cite{iso26300}.
ODT는 Microsoft Word(이하 Word), 한컴 한/글(이하 한/글), 오픈소스 오피스 프로그램 LibreOffice~\cite{libreoffice} 등 주요 워드프로세서에 활용할 수 있어 특정 제품에 종속되지 않는다. 따라서 제안하는 평가를 워드프로세서 문서 복원 전반에 적용할 수 있다

% 출력본 → 복원 방식 → 복원본 → 편집 → 완성본의 개념도.
% 원본은 PDF 생성·편집 시간 비교에 사용하며, 준비 과정과 시간 정규화는 본문에서 설명한다.
% 편집 화살표는 평가 목표를 나타내며 실제 편집 실행을 뜻하지 않는다.
\begin{figure}[H]
\centering
\begin{tikzpicture}[
  docname/.style={font=\footnotesize,align=center,anchor=north,inner sep=1pt},
  proc/.style={draw=black!70,rounded corners=2pt,font=\footnotesize,minimum width=14mm,minimum height=6mm,inner sep=2pt},
  metric/.style={font=\footnotesize\bfseries,anchor=north,inner sep=1pt},
  span/.style={decorate,decoration={brace,mirror,amplitude=3pt},black!60}]
% 세 문서는 같은 크기의 표로 그리고, 완성본만 값 칸을 채운다.
\newcommand{\miniform}[2]{%
  \begin{scope}[shift={(#1,0)}]
    \draw[black!70] (-0.65,0.30) rectangle (0.65,-0.30);
    \draw[black!70] (-0.20,0.30) -- (-0.20,-0.30) (-0.65,0) -- (0.65,0);
    \foreach \y in {-0.15,0.15} {
      \draw[black!50,line width=0.8pt] (-0.54,\y) -- (-0.31,\y);
      \ifnum#2=1
        \draw[black,line width=1pt] (-0.07,\y) -- (0.52,\y);
      \fi
    }
  \end{scope}}
% 과정: 문서 중심 x = 0.65, 4.35, 8.05 / 처리 단계 중심 x = 2.50, 6.20
\miniform{0.65}{0}
\miniform{4.35}{0}
\miniform{8.05}{1}
\node[proc] (restore) at (2.50,0) {복원 방식};
\node[proc] (edit) at (6.20,0) {편집};
\draw[paperarrow] (1.30,0) -- (restore.west);
\draw[paperarrow] (restore.east) -- (3.70,0);
\draw[paperarrow] (5.00,0) -- (edit.west);
\draw[paperarrow] (edit.east) -- (7.40,0);
\node[docname] at (0.65,-0.36) {출력본(PDF)};
\node[docname] at (4.35,-0.36) {복원본(ODT)};
\node[docname] at (8.05,-0.36) {완성본(ODT)};
% 측정: 시각 유사도는 출력본과 복원본의 외관 비교, 편집 시간은 복원본을 완성본으로 편집하는 시간
\draw[span] (0.65,-0.80) -- (4.35,-0.80);
\draw[span] (4.35,-0.80) -- (8.05,-0.80);
\node[metric] at (2.50,-0.94) {시각 유사도};
\node[metric] at (6.20,-0.94) {편집 시간};
\end{tikzpicture}
\caption{평가 과정}
\label{fig:pipeline}
\end{figure}


\subsection{시각 유사도}

시각 유사도는 복원본의 외관이 출력본과 얼마나 유사한지를 나타낸다.
문서를 이미지로 그리는 과정을 \textbf{렌더링}이라 한다.
출력본과 복원본을 렌더링한 뒤, 같은 페이지 번호끼리 대응시켜 이미지 유사도 지표인 SSIM~\cite{wang2004ssim}으로 비교한다(최대 1, 클수록 유사).
글자 경계의 미세한 그리기 차이는 비교 전에 가우시안 블러로 완화한다.
문서 전체의 값은 출력본을 기준으로 대응 이미지의 점수를 평균하며, 복원본에서 누락된 이미지는 0으로 계산한다.
추가된 페이지는 이 평균에 반영되지 않으므로, 출력본보다 페이지 수가 늘어난 복원본의 비율을 \textbf{페이지 넘침}으로 따로 보고한다.

\subsection{편집 시간}

편집 시간은 복원본을 완성본으로 만드는 데 필요한 시간의 추정값으로, 낮을수록 편집이 용이한 것으로 본다.
문서를 문단, 표, 글상자, 이미지 등의 \textbf{블록}으로 나누고 필요한 작업을 블록의 추가, 변경, 삭제로 나타낸다.
예를 들어 표의 셀 텍스트를 바꾸거나 행을 추가할 수 있으며, 기존 표를 삭제하고 다시 만드는 경로도 허용한다.
정의한 블록 대응에서 시간이 적게 드는 편집 조합을 동적 계획법으로 구한다.
목표는 내용과 표 구조가 완성본과 같고, 대응 글자의 위치 차이가 같은 페이지에서 가로·세로 12pt(약 4.2mm) 이내인 문서다.

각 편집의 시간은 숙련 사용자의 동작 시간을 합산하는 KLM(Keystroke-Level Model)~\cite{card1980klm}으로 추정한다.
키 입력 0.28초, 포인터 이동 1.1초, 클릭 0.2초 등의 시간을 가정하고, 문자 종류와 입력기에 따라 문자당 키 수를 달리 적용한다.
글꼴·정렬·간격 등 글자·문단 모양의 변경도 필요한 편집에 포함한다.
이 값들은 실험의 모델 설정이며 사용자 실측 시간이 아니다.
문서의 분량과 수정량에 따른 차이를 줄이기 위해 복원본의 편집 시간을 원본에서 같은 완성본을 만드는 시간으로 나눈 배수로 보고한다.
1배는 원본과 같은 시간, 2배는 두 배의 추정 시간을 뜻한다.

\section{실험 설정}

원본 321개와 각 원본에 대응하는 완성본을 동일 출처에서 수집하고 모두 수작업으로 검수했다.
수집 당시 원본의 파일 형식은 DOCX 235개, ODT 48개, HWPX 38개이며, 각 원본과 해당 완성본의 형식은 같다.
원본의 언어는 일본어 195개, 중국어 50개, 한국어와 영어 각 37개 등이고, 출처 웹 도메인은 189곳이다.

DOCX와 HWPX는 각각 Word와 한/글의 기본 형식이어서, 같은 문서도 다른 프로그램에서 열면 배치가 달라질 수 있다.
공통 형식에서 평가하도록 DOCX와 HWPX의 원본과 완성본은 평가 전에 ODT로 변환했다.
렌더링 프로그램에 따른 차이를 줄이기 위해 원본과 복원본은 모두 LibreOffice로 렌더링한다.
원본을 렌더링한 PDF를 출력본으로 사용하며, 이는 복원 방식의 입력이자 시각 유사도의 비교 기준이다.
시각 유사도는 72 dpi 이미지에 표준편차 1픽셀의 가우시안 블러를 적용해 계산한다.

복원 방식은 대표적인 세 계열에서 네 가지를 선택했다.
워드프로세서의 PDF 가져오기 계열로 Word와 한/글, OCR 기반 변환 도구 계열로 글자와 영역·표 구조 인식을 결합한 Docling, 시각언어모델(VLM) 계열로 GPT-6-Luna를 사용한다.

대조군으로는 \textbf{원본} 자체를 사용하는 경우와 출력본을 렌더링해 삽입하는 \textbf{이미지} 방식을 함께 평가한다.
전자는 결과의 오차를 점검하고, 후자는 외관만 재현한 경우의 편집 시간을 보여준다.

\section{실험 결과}

\subsection{복원 방식별 비교}

표~\ref{tab:main-results}는 복원 방식별 시각 유사도와 편집 시간의 중앙값, 페이지 넘침 비율을 보인다.
원본에서는 셀 몇 개만 기입하면 되더라도 복원본에서는 표와 내용을 다시 만들어야 하는 경우가 있어, 편집 시간이 원본의 수십 배로 커질 수 있다(최대 약 34배).
이러한 극단값의 영향을 줄이기 위해 평균 대신 중앙값으로 비교한다.
% Source: reports/experiments/main/rows.csv. 최대 33.7배, KLM 시간 중앙값과 배수 중앙값은 서로 다른 집계.
원본을 복원 결과로 사용했을 때 시각 유사도는 모두 1.00, 편집 시간은 모두 1.00배였다.
원본의 편집 시간 중앙값은 236초(약 4분)다.
네 방식 중 가장 짧은 Word도 2.34배로 원본의 두 배를 넘었고 Docling은 2.77배였다.
페이지 넘침은 GPT-6-Luna가 34.9\%로 가장 높고 한/글이 17.1\%로 가장 낮았다.

% Source: reports/experiments/main/summary.json (실험 current: 원본과 완성본을 ODT로 전처리, docs/experiment.md)
% created_at: 2026-10-18T22:09:46Z; run_id: main; unit: seconds; tasks: 321
% GPT-6-Luna = png_150__gpt6luna_odt, Docling = vector_pdf__docling_odt (ODT 직접 출력); MinerU2.5는 논문 비교 대상에서 제외.
% Columns: labels.*.{ssim_blur_native_pages_median,ratio_median,extra_pages_share}
% 시각 유사도·편집 시간(ratio, 원본 대비 배수)은 개별 평가값의 중앙값(소수 두 자리), 페이지 넘침은 출력본보다 페이지 수가 늘어난 복원본의 비율(백분율 한 자리).
% 편집 시간은 원본 321개에 대해 실패를 포함한 결과; 시각 유사도와 페이지 넘침은 시각 평가가 있는 성공 건 기준.
% 양성 대조군 먼저, 실제 방식은 ratio_median 오름차순, 음성 대조군 마지막.
\begin{table}[!htb]
\centering
\caption{복원 방식별 평가 결과}
\label{tab:main-results}
\papertablestyle
\begin{tabularx}{\linewidth}{@{}Lrrr@{}}
\toprule
방식 & \begin{tabular}[b]{@{}r@{}}시각 유사도\\(최대 1)\end{tabular} & \begin{tabular}[b]{@{}r@{}}편집 시간\\(배)\end{tabular} & \begin{tabular}[b]{@{}r@{}}페이지 넘침\\(\%)\end{tabular} \\
\midrule
원본 & 1.00 & 1.00 & 0.0 \\
\midrule
Word & 0.73 & \textbf{2.34} & 26.8 \\
한/글 & 0.71 & 2.49 & 17.1 \\
GPT-6-Luna & 0.73 & 2.67 & 34.9 \\
Docling & 0.71 & 2.77 & 20.9 \\
\midrule
이미지 & 0.99 & 3.51 & 0.0 \\
\bottomrule
\end{tabularx}
\end{table}


추가 작업의 구성을 살펴보면, 네 방식은 표·문단 등의 블록을 새로 넣는 데 전체 추정 편집 시간의 54--64\%를 사용했다(그림~\ref{fig:edit-breakdown}).
이는 각 방식의 추정 편집 시간을 모두 합산한 뒤 구한 비중이다.
추가는 누락된 블록을 보완하는 작업뿐 아니라 기존 블록을 삭제하고 다시 만드는 작업도 포함한다.
반면 원본은 기존 블록의 변경이 67\%를 차지했다.
따라서 추가 비중을 요소별 복원 정확도로 해석하지 않고, 수정 목표에 이르기 위해 다시 만들어야 하는 작업의 비중으로 해석한다.
그림의 기타는 삭제와 페이지 설정을 합한 값이다.
% Source: reports/experiments/main/edit_tasks.json labels.*.by_operation_share_pct.
% add: Word 55.06, 한/글 56.91, GPT-6-Luna 54.45, Docling 64.47; 원본 edit 67.17.

% Source: figures/edit-breakdown.dat (figures/edit-breakdown.py가 reports/experiments/main/edit_tasks.json에서 생성)
% 원본 321개에 대한 전체 편집 시간를 변경, 추가, 삭제·페이지 설정로 나눈 비율(%).
% 막대는 원자료의 소수 두 자리를 유지하고, 추가 비중은 여기에서 한 번만 정수 반올림한다.
\begin{figure}[H]
\centering
\begin{tikzpicture}
\begin{axis}[
  xbar stacked, bar width=3.1mm, scale only axis, width=0.76\linewidth, height=24mm,
  xmin=0, xmax=100, xtick={0,25,50,75,100}, xlabel={시간 비중(\%)},
  y dir=reverse, ytick=data, yticklabels from table={figures/edit-breakdown.dat}{name},
  tick label style={font=\scriptsize}, label style={font=\footnotesize}, xlabel shift=-2pt,
  enlarge y limits=0.08, axis line style={black!60}, tick style={black!60},
  legend style={font=\scriptsize, at={(0.5,1.02)}, anchor=south, legend columns=3, draw=none,
    /tikz/every even column/.append style={column sep=6pt}},
  legend image code/.code={\draw[#1] (0cm,-0.1cm) rectangle (0.3cm,0.1cm);}]
\addplot[fill=black!22, draw=black!60] table[x=edit, y expr=\coordindex] {figures/edit-breakdown.dat};
% 추가 구간의 중앙에 해당 비중만 표시한다.
\addplot[fill=black!65, draw=black!65,
  nodes near coords={\pgfmathprintnumber[fixed,precision=0]{\pgfplotspointmeta}},
  every node near coord/.append style={font=\scriptsize,text=white}]
  table[x=add, y expr=\coordindex] {figures/edit-breakdown.dat};
\addplot[fill=white, draw=black!60] table[x=del, y expr=\coordindex] {figures/edit-breakdown.dat};
\legend{변경, 추가, 기타}
\end{axis}
\end{tikzpicture}
\caption{편집 유형별 시간 비중}
\label{fig:edit-breakdown}
\end{figure}


위치 허용 범위를 6pt와 24pt로 바꾸거나, 키 입력 시간을 절반과 두 배로 바꾸거나, 문자당 키 수를 모두 1로 두어도 방식별 편집 시간 중앙값의 순서는 유지되었다.
% Source: reports/experiments/main/sensitivity.json, 다섯 조건 Kendall tau 1.0.

\balance
\subsection{시각 유사도와 편집 시간의 관계}

그림~\ref{fig:ssim-ratio}의 각 점은 출력본 하나를 특정 방식으로 복원했을 때의 시각 유사도와 편집 시간을 나타낸다.
세로축은 같은 배수 변화가 같은 간격이 되는 로그 눈금이며, 점선은 원본과 같은 편집 시간인 1배를 나타낸다.
대조군을 제외한 네 방식의 스피어만 순위 상관계수는 $-0.24$로, 시각 유사도가 높을수록 편집 시간이 짧아지는 관계는 약했다.
출력본과 외관이 거의 같은 시각 유사도 0.9 이상의 54건에서도 편집 시간 중앙값은 원본의 1.98배였다.
이는 외관의 유사성만으로 편집 용이성을 판단하기 어렵다는 점을 보여준다.

% Source: figures/ssim-ratio.dat (figures/ssim-ratio.py가 reports/experiments/main/rows.csv에서 생성)
% 점: 대조군과 MinerU2.5를 제외한 네 방식의 평가 결과 1,267건 (ssim_blur_native_pages, ratio).
% 대조군 표시는 summary.json labels.{png_150__image_only,oracle__template_odt}의
%   ssim_blur_native_pages_median, ratio_median (0.99, 3.51), (1.00, 1.00).
\begin{figure}[H]
\centering
\begin{tikzpicture}
\begin{axis}[
  width=\linewidth, height=40mm,
  xmin=0.1, xmax=1.02, ymode=log, ymin=0.4, ymax=40,
  ytick={0.5,1,2,5,10,20}, yticklabels={0.5,1,2,5,10,20},
  xtick={0.2,0.4,0.6,0.8,1.0},
  tick label style={font=\scriptsize}, label style={font=\footnotesize},
  xlabel={시각 유사도(최대 1)}, ylabel={편집 시간(배)},
  ylabel shift=-3pt, xlabel shift=-2pt,
  axis line style={black!60}, tick style={black!60},
  clip marker paths=true]
\addplot[only marks, mark=*, mark size=0.7pt, draw opacity=0, fill=black, fill opacity=0.28]
  table[x=ssim, y=ratio] {figures/ssim-ratio.dat};
\addplot[black!50, dashed, domain=0.1:1.02, samples=2] {1};
\addplot[only marks, mark=square*, mark size=2.2pt, fill=white, draw=black, semithick] coordinates {(0.9936,3.508)};
\node[font=\scriptsize, anchor=east, fill=white, inner sep=1pt] at (axis cs:0.965,6.2) {이미지};
\draw[black!70] (axis cs:0.965,6.2) -- (axis cs:0.99,3.6);
\addplot[only marks, mark=triangle*, mark size=2.4pt, fill=white, draw=black, semithick] coordinates {(1.0,1.0)};
\node[font=\scriptsize, anchor=north east, fill=white, inner sep=1pt] at (axis cs:0.90,0.82) {원본};
\end{axis}
\end{tikzpicture}
\caption{시각 유사도와 편집 시간의 관계}
\label{fig:ssim-ratio}
\end{figure}


\section{결 론}

본 논문의 편집 시간은 모델 추정값으로, 사용자의 실제 작업 시간과 비교하지 않았다는 한계가 있다.
그러나 네 가지 복원 방식에 적용한 결과, 시각 유사도만으로는 드러나지 않는 수정 작업의 차이를 정량적으로 비교할 수 있음을 보였다.
이는 PDF 양식 복원을 편집 용이성이라는 관점에서 평가하는 방법을 제시했다는 점에서 의의가 있다.
향후 더 다양한 양식에 대한 평가와 사용자 실험으로 지표의 타당성을 검증한다면, 실제 수정 부담을 줄이는 복원 방식을 개발하고 평가하는 데 활용할 수 있을 것이다.

{\small
\bibliographystyle{ieeetr}
\bibliography{ref}
}
\end{document}
